% Generated by scripts/build_blueprint_tex.py from data/model.json -- do not edit by hand.
\begin{bptimelinetable}
\bptlrow{1}{10/6/26}{6 October 2026}{Stepanov determinant}{3}{Historical motivation}{\bpcompref{family}}{The single-line entry records an origin of the argument. Its underlying function-field determinant is unspecified.}{stratum-est}
\bptlrow{2}{10/9/26}{9 October 2026}{Invariant differentials on a CM curve}{3}{Historical calculation}{\bpcompref{family}}{The entry describes inert-prime divisibility of iterated derivatives. The argument for this displayed exponent was not identified in the selected source path.}{stratum-est}
\bptlrow{3}{10/19/26+}{19 October 2026 onwards}{The finite-field sign argument}{3}{Historical analogy}{\bpnocomp}{The footnote compares a zero mode with an error bounded by iterated Parseval. Its objects and hypotheses differ from the Dirichlet L-function endpoint.}{tier-motivation}
\bptlrow{4}{10/6/26+}{6 October 2026 onwards}{Determinantal positivity with multiplicity one}{3}{Historical motivation}{\bpnocomp}{The short entry leaves W and S implicit. An exact theorem-to-type identification would require definitions of both.}{tier-motivation}
\bptlrow{5}{10/20/26}{20 October 2026}{The smoothed von Mangoldt sum}{3}{Formula matched}{\bpcompref{statistic}}{The displayed zeta error has the exact ramp and the correction (1 + Euler’s constant)x. Character twists enter the level average. The sentence continues onto page 4.}{stratum-ana}
\bptlrow{6}{10/8/26}{8 October 2026}{Two-support positivity test}{4}{Historical calculation}{\bpcompref{criterion}}{The two indicator functions are a proposed model for a zero-free exponent. The final power-window criterion has a different statement.}{stratum-ana}
\bptlrow{7}{10/14/26}{14 October 2026}{Interlacing weights and periods}{4}{Supporting comparison}{\bpcompref{inputs}}{A local Bessel model is constructed from two interlacing conditions. The announcement’s claim about determinant values 0 and 1 has further content.}{stratum-arith}
\bptlrow{8}{11/15/26}{15 November 2026}{A chain with square-root length}{4}{Quantitative comparison}{\bpcompref{family}}{The arithmetic rank specification puts m between one and two times x raised to 1/2 minus epsilon/16. The finite kernel recursion explains iteration.}{stratum-est}
\bptlrow{9}{11/15/26}{15 November 2026}{Norms, levels and arithmetic periods}{4}{Supporting comparison}{\bpcompref{family}}{One result relates a CM period to a logarithmic derivative. Another cancels a level volume against an index. The complete order-m cancellation requires more data.}{stratum-est}
\bptlrow{10}{11/15/26}{15 November 2026}{Divisibility of repeated lifts}{4}{Supporting comparison}{\bpcompref{family}}{A specific minor vanishes modulo a divisor and a lifted column has a higher-order remainder. Turning these local facts into the required global gain is an additional step.}{stratum-est}
\bptlrow{11}{10/6/26}{6 October 2026}{The tower of fields}{4}{Supporting comparison}{\bpcompref{cohomology}}{The formal window ranges over fields specified by a fixing subgroup in a cyclotomic extension and by layer parameter e(x).}{stratum-arith}
\bptlrow{12}{11/15/26}{15 November 2026}{Seesaw and Bessel periods}{4}{Supporting comparison}{\bpcompref{inputs}}{Conditional unfolding identifies a theta-lift period with an integral. A middle-step volume identity is independent of the admissible character. A full iterated comparison is a stronger assertion.}{stratum-arith}
\bptlrow{13}{10/6/26}{6 October 2026}{The slope inequality}{4}{Mathematical statement matched}{\bpcompref{family}}{The Euclidean lattice model bounds the slope of an injectively mapped source by the maximal target slope plus the logarithm of the map bound. Identification of the particular period lattice is separate.}{stratum-est}
\bptlrow{14}{11/15/26}{15 November 2026}{The provisional exponent two-thirds}{4}{Intermediate strategy}{\bpcompref{comparison}}{The author describes a barrier for an earlier estimate. The endpoint reads a one-half-plus-loss exponent. The earlier barrier is not derived in this comparison.}{stratum-est}
\bptlrow{15}{11/15/26}{15 November 2026}{A common integral structure across characters}{4}{Supporting comparison}{\bpcompref{family}}{The final deduction imposes local and global controls on the same family. A precise comparison with the gcd in the historical description remains open.}{stratum-est}
\bptlrow{16}{10/14/26}{14 October 2026}{Period pairings and approximate orthogonality}{4}{Supporting comparison}{\bpcompref{comparison}}{Unfolding relates a Poincare-series pairing to an integral of periods. The later host comparison turns suitable bounds into an estimate for the weighted prime sum.}{stratum-est}
\bptlrow{17}{11/15/26}{15 November 2026}{Cohomology and a weakening of Vandiver}{4}{Quantitative comparison}{\bpcompref{cohomology}}{The window requires finite lengths and an aggregate weighted bound. Further implications for individual K-groups are conjectural comments in the announcement.}{stratum-arith}
\bptlrow{18}{11/15/26}{15 November 2026}{Covolumes, lattice indices and congruences}{4}{Supporting comparison}{\bpcompref{inputs}}{A covolume bound is an explicit premise of the automorphic assembly. The local congruence lemma gives a representative higher-order equality.}{stratum-arith}
\bptlrow{19}{10/19/26}{19 October 2026}{A Behrend set in the tower construction}{4}{Supporting comparison}{\bpcompref{primes}}{Encoding a populated integer sphere in a large base excludes short affine relations and controls intersections with short orbits.}{stratum-arith}
\bptlrow{20}{11/16/26}{16 November 2026}{Dirichlet GRH}{4}{Endpoint statement matched}{\bpcompref{endpoints}}{The source declares that zeros of every Dirichlet L-function in the open critical strip are on the critical line. This entry records source-statement correspondence.}{stratum-dest}
\end{bptimelinetable}
